\section{Pacuit's canonical neighborhood constructions}
\label{sec:p08-n-b-006-04}
\label{sec:p08-pacuit-neighborhood-source}

Pacuit's published 2017 \emph{Core Theory} chapter develops extensional
modal systems through proof sets and a choice of canonical neighborhood
functions. Minimal canonical neighborhoods suffice for E, EC and EN;
monotonic systems require a separate supplementation argument. Filters
and augmented neighborhoods give two different semantic presentations
of K \cite[\S\S2.2.1, 2.3, 2.5, pp.~47--49, 52--66, 80--82]{Pacuit2017}.
This treatment retains those constructions and their local-consequence
scope. It does not infer that every neighborhood logic is complete for
unrestricted ordinary frames merely because some canonical model exists.

\subsection{The source vocabulary and exact modal operations}

\begin{hbdefinition}{def:p08-pn-language-model}{Language, models and truth sets}
Let \(\mathsf{At}\) be finite or countable and let
\(\mathcal L(\mathsf{At})\) be its classical propositional modal
language. Use \(\Box\) as primitive and
\(\Diamond\varphi\defeq\neg\Box\neg\varphi\); the two-operator
notation in the source's Dual schema is interpreted through this
specified definitional expansion. Other Boolean connectives are
abbreviations for negation and conjunction. When \(\mathsf{At}\)
is empty, include \(\bottom,\topv\) explicitly; otherwise they
can be defined by a contradiction and its negation.

A frame is \(\mathcal F=\langle W,N\rangle\), with nonempty W
and \(N:W\to\wp(\wp(W))\). Its model is
\(\mathcal M=\langle W,N,V\rangle\), with
\(V:\mathsf{At}\to\wp(W)\). Here \(\wp\) is the source
powerset notation and
\([![\varphi]!]_{\mathcal M}=\{w:\mathcal M,w\models\varphi\}\)
is its truth set. Boolean truth is classical, and
\[
\begin{aligned}
 \mathcal M,w\models\Box\varphi
   &\ \Longleftrightarrow\ [![\varphi]!]_{\mathcal M}\in N(w),\\
 \mathcal M,w\models\Diamond\varphi
   &\ \Longleftrightarrow\ W\setminus[![\varphi]!]_{\mathcal M}
                         \notin N(w).
\end{aligned}
\]
The alternative source operator \(\langle\ ]\varphi\) means
that some \(X\in N(w)\) is a subset of its truth set. It is not
the ordinary diamond and is not the same as exact-membership box
on arbitrary neighborhoods.
\end{hbdefinition}

\begin{hblemma}{thm:p08-pn-extensional}{proved=annex:p8.pn.extensional}{Extensionality, duality and the alternative monotonic reading}
Equal truth sets have equal boxed truth sets. Validity of
\(\varphi\leftrightarrow\psi\) throughout a model therefore
implies validity of \(\Box\varphi\leftrightarrow\Box\psi\)
throughout that model; merely equal truth values at one world do not.
The Dual schema \(\Box\varphi\leftrightarrow\neg\Diamond\neg\varphi\)
is valid. The alternative operator \(\langle\ ]\) is monotonic
on every neighborhood model and agrees with \(\Box\) on supplemented
models. Box itself need not be monotonic on an arbitrary model.
This supplies the semantic bridge to the existing defined-diamond sheets
without treating an undeclared additional primitive as an abbreviation.
\end{hblemma}

\begin{hbdefinition}{def:p08-pn-closure-properties}{Independent neighborhood closure conditions}
A family \(\mathcal U\subseteq\wp(W)\) is supplemented, or
monotonic, when it is closed under supersets. It is intersection-closed
when \(X,Y\in\mathcal U\) implies \(X\cap Y\in\mathcal U\).
Equivalently it is closed under \emph{nonempty} finite intersections;
the empty intersection would instead impose the independent unit
condition \(W\in\mathcal U\).
A filter has supplementation, binary intersection closure and the unit.
A proper filter also excludes \(\varnothing\). Regular neighborhoods
need only supplementation and binary intersections; their family may
be empty. A family is augmented when it is supplemented and contains
its core \(\bigcap\mathcal U\), with the empty-family core taken
relative to W. An augmented family is nonempty and contains the unit.
These properties are imposed separately at every world of a frame.
Pacuit also calls an arbitrary neighborhood family \emph{proper} when
it never contains both a set and its complement. This source use of
``proper'' is distinct from the empty-set exclusion in the definition
of a proper filter; their equivalence needs the regularity hypotheses
proved below.
\end{hbdefinition}

\begin{hbtheorem}{thm:p08-pn-correspondence}{proved=annex:p8.pn.correspondence}{All-valuations M, C, N and properness correspondences}
For the arbitrary-set converse directions assume at least two available
atoms. The frame correspondences are exactly
\[
\begin{array}{lll}
 \mathrm M:&\Box(\varphi\land\psi)\to(\Box\varphi\land\Box\psi)
      &\text{supplementation},\\
 \mathrm C:&(\Box\varphi\land\Box\psi)\to\Box(\varphi\land\psi)
      &\text{binary intersection closure},\\
 \mathrm N:&\Box\topv&\text{unit containment}.
\end{array}
\]
Each condition implies its schema's validity for any alphabet.
With one atom available, validity of
\(\Box\varphi\to\Diamond\varphi\) corresponds to absence
of complementary pairs \(X,W\setminus X\) in any \(N(w)\).
Validity of \(\Diamond\topv\) instead says
\(\varnothing\notin N(w)\). The latter two properties agree
on regular families but not on all monotonic families
\cite[lemmas~2.19--2.22]{Pacuit2017}.
\end{hbtheorem}

\subsection{Theorem systems, naming and premise discipline}

\begin{hbdefinition}{def:p08-pn-logics}{E and the selected schema extensions}
Pacuit's E is the least theorem set containing propositional tautologies
and Dual, closed under modus ponens and theorem-only replacement of
equivalents. EM, EC and EN add M, C and N respectively; EMC adds M
and C; EMCN adds all three. The source's ECM names the same M/C
extension as EMC. These names are not the standalone modal operators
or conditions M, C and N. The registered sheets
\ref{sheet:hilbert-e}, \ref{sheet:hilbert-m}, \ref{sheet:hilbert-c},
\ref{sheet:hilbert-n} and \ref{sheet:hilbert-r} give their existing
rule surfaces in the defined-diamond language.

The source calls E the least minimal or classical modal logic, EM the
least monotonic logic and EMC the least regular logic. EN is not normal
merely because necessitation is admissible; K additionally requires its
distribution schema. EK means E plus that schema, without adding
necessitation, and is different from K.
\end{hbdefinition}

\begin{hbtheorem}{thm:p08-pn-derived-rules}{proved=annex:p8.pn.derived-rules}{Source replacement, RM, necessitation and the normal boundary}
Theorem equivalence may be replaced at any chosen occurrences in a
formula, including under its modalities. EM is exactly E closed
under theorem-only monotonicity of box on theorem implications.
EN is exactly E closed under theorem-only necessitation.
EMC proves the distribution schema K without acquiring N;
EMCN is exactly the normal logic K. These identities preserve
theorems and the finite-implication local deduction relation below;
they do not permit either modal rule on an undischarged local premise.
More generally, a classical L is regular exactly when it preserves
boxed theorem consequences with every positive finite number of
antecedents; it is normal when the zero-antecedent case is included
as well. The empty case is necessitation, not an innocuous convention.
\end{hbtheorem}

\begin{hbtheorem}{thm:p08-pn-derived-principles}{proved=annex:p8.pn.derived-principles}{The source's monotonic and regular consequence package}
E proves \(\Diamond\topv\leftrightarrow\neg\Box\bottom\)
and \(\Box\topv\leftrightarrow\neg\Diamond\bottom\).
Every monotonic logic proves
\[
\begin{aligned}
 \Box\psi&\to\Box(\varphi\to\psi),&
 \Box\neg\varphi&\to\Box(\varphi\to\psi),\\
 \Diamond(\varphi\to\psi)&\lor\Box(\psi\to\varphi),&
 \Diamond(\varphi\land\psi)&\to(\Diamond\varphi\land\Diamond\psi).
\end{aligned}
\]
Theorem implication also gives theorem diamond monotonicity.
Every regular logic proves
\[
\begin{aligned}
 \Box(\varphi\to\psi)&\to(\Diamond\varphi\to\Diamond\psi),\\
 \Box(\varphi\leftrightarrow\psi)&\to(\Box\varphi\leftrightarrow\Box\psi),\\
 (\Box\varphi\land\Diamond\psi)&\to\Diamond(\varphi\land\psi).
\end{aligned}
\]
Moreover \((\Box\varphi\to\Diamond\varphi)\to\Diamond\topv\)
is monotonic, while its converse requires regularity; hence
\(\Diamond\topv\leftrightarrow(\Box\varphi\to\Diamond\varphi)\)
is regular. A monotonic countermodel to the converse was specified
in the correspondence proof \cite[ex.~2.17]{Pacuit2017}.
\end{hbtheorem}

\begin{hbtheorem}{thm:p08-pn-finite-separations}{proved=annex:p8.pn.finite-separations}{Complete separating models and the EK boundary}
With two available atoms, finite neighborhood frames separate M, C and N as independent
conditions. EN can fail distribution, and EMC can fail N.
Pacuit's four-world EK example validates distribution under every
valuation but fails both M and C. Thus distribution without
necessitation does not characterize either regularity or normality.
On a unit frame, \(\{p\}\not\models^{\ell}\Box p\), even
though global \emph{model} truth of p implies model truth of
\(\Box p\). The reciprocal proof specifies every neighborhood
family and valuation, including the previously unspecified worlds
of the source's root-only counterexample.
\end{hbtheorem}

\begin{hbdefinition}{def:p08-pn-local-deduction}{Source local consequence and strong completeness}
For a logic L and any formula set \(\Gamma\), Pacuit's
\(\Gamma\vdash_L\varphi\) means that some finite
\(\alpha_1,\ldots,\alpha_k\in\Gamma\) have
\(\vdash_L(\alpha_1\land\cdots\land\alpha_k)\to\varphi\).
For k=0 use theoremhood of \(\varphi\). This is equivalent
to finite premise derivations on the sheets with theorem-only modal
inferences. Consistency means \(\Gamma\not\vdash_L\bottom\).
Local semantic consequence \(\Gamma\models^{\ell}_{\mathscr F}\varphi\)
requires every world in every model on \(\mathscr F\) satisfying
all premises to satisfy \(\varphi\). Strong completeness permits
arbitrary, not necessarily finite or effectively specified, premise sets.

Model-global consequence assumes the premises at every world of one
model. Frame-global consequence instead assumes their validity under
every valuation on one frame. These are different relations. In
particular a nonempty frame never validates an unconstrained atom p,
so a frame-validity premise p makes its consequence vacuous. Neither
global convention is substituted for the source's local completeness
relation \cite[defs.~2.30--2.34]{Pacuit2017}.
\end{hbdefinition}

\subsection{Proof sets and the two canonical extrema}

\begin{hbdefinition}{def:p08-pn-proof-sets}{Maximal theories and proof sets}
For consistent L, let \(M_L\) be its maximal consistent formula sets.
Equivalently they are consistent sets deciding every formula by
including it or its negation. Write
\(|\varphi|_L=\{\Delta\in M_L:\varphi\in\Delta\}\)
for the source proof set, and
\(P_L=\{|\varphi|_L:\varphi\in\mathcal L(\mathsf{At})\}\)
for the family of all proof sets. A proof set is a subset of the
canonical world set, not a set of derivations or a numerical size.
Countability of the language gives at most countably many proof sets;
no unnecessary universal cardinality assumption on \(M_L\) is used.
\end{hbdefinition}

\begin{hblemma}{thm:p08-pn-proof-set-lemmas}{proved=annex:p8.pn.proof-set-lemmas}{The complete Lindenbaum and proof-set interface}
Every consistent premise set has a maximal consistent extension.
Maximal theories contain all theorems and have every classical Boolean
membership property. Their proof sets satisfy
\[
\begin{aligned}
 |\neg\varphi|_L&=M_L\setminus|\varphi|_L,\\
 |\varphi\land\psi|_L&=|\varphi|_L\cap|\psi|_L,\\
 |\varphi|_L\subseteq|\psi|_L
   &\ \Longleftrightarrow\ \vdash_L\varphi\to\psi,\\
 |\varphi|_L=|\psi|_L
   &\ \Longleftrightarrow\ \vdash_L\varphi\leftrightarrow\psi.
\end{aligned}
\]
Also \(\psi\in\bigcap|\varphi|_L\) exactly when
\(\vdash_L\varphi\to\psi\), with the empty family intersection
taken in the formula universe. These results hold for finite or
countable alphabets, including arbitrary infinite initial premises.
In particular \(|\topv|_L=M_L\), \(|\bottom|_L=\varnothing\)
and \(|\varphi\lor\psi|_L=|\varphi|_L\cup|\psi|_L\).
\end{hblemma}

\begin{hbdefinition}{def:p08-pn-canonical}{Minimal and maximal canonical neighborhoods}
The canonical valuation is \(V_L(p)=|p|_L\). A neighborhood
function on \(M_L\) is canonical for L when
\(|\varphi|_L\in N_L(\Gamma)\) exactly when
\(\Box\varphi\in\Gamma\). The source's extrema are
\[
\begin{aligned}
 N_L^{\min}(\Gamma)&=\{|\varphi|_L:\Box\varphi\in\Gamma\},\\
 N_L^{\max}(\Gamma)&=N_L^{\min}(\Gamma)
                 \cup(\wp(M_L)\setminus P_L).
\end{aligned}
\]
They share the worlds and valuation but may differ on non-proof sets.
It must be proved that the first expression is canonical: equal proof
sets must not have incompatible boxed membership. A specified canonical
condition alone does not establish existence of its function.
\end{hbdefinition}

\begin{hbtheorem}{thm:p08-pn-canonical-truth}{proved=annex:p8.pn.canonical-truth}{Existence, extrema and every truth-lemma case}
The two displayed functions are canonical for every consistent
classical L with theorem-only replacement. Every canonical function
lies between them pointwise. In any corresponding canonical model,
\([![\varphi]!]_{\mathcal M_L}=|\varphi|_L\) for every
formula, including inconsistent formulas. Consequently every consistent
premise set is true at one canonical world. This model-truth assertion
does not by itself prove that an arbitrary valuation on that same
frame validates L.
\end{hbtheorem}

\begin{hbtheorem}{thm:p08-pn-minimal-not-monotone}{proved=annex:p8.pn.minimal-not-monotone}{Why the minimal EM canonical function needs supplementation}
For EM there is a world \(\Gamma\), a proof set X in
\(N^{\min}_{EM}(\Gamma)\), and a non-proof superset Y of X
that is absent from that family. Hence the minimal canonical
function is not necessarily supplemented. The proof supplies an
actual non-proof subset disjoint from X before forming Y; it does
not infer that an arbitrary non-proof set remains non-proof after
union. The construction also covers a language with no atoms by
using its stipulated constants and modal formulas.
\end{hbtheorem}

\subsection{Supplementation and all six exact completeness applications}

\begin{hbdefinition}{def:p08-pn-supplementation}{The source supplementation operation}
For any family \(\mathcal U\subseteq\wp(W)\), let
\(\mathcal U^{\rm mon}=\{X\subseteq W:\exists Y\in\mathcal U,
Y\subseteq X\}\). Supplement a model by applying this operation
to every \(N(w)\), leaving its worlds and atomic valuation fixed.
For a canonical model this closure must be checked not to insert
a previously rejected \emph{proof set}; arbitrary new non-proof
neighborhoods may otherwise be harmless to canonical truth.
\end{hbdefinition}

\begin{hbtheorem}{thm:p08-pn-supplemented-canonical}{proved=annex:p8.pn.supplemented-canonical}{Canonical supplementation and preservation of C and N}
When L contains M, supplementation of its minimal canonical
function remains canonical. When L also contains C, the supplemented
families are binary intersection-closed; when L contains N, they
contain \(M_L\). Independently, a minimal canonical family is
binary intersection-closed when L contains C and contains the unit
when L contains N. Thus the required properties hold together in
the same canonical model, not merely in separately chosen models.
\end{hbtheorem}

\begin{hbtheorem}{thm:p08-pn-six-completeness}{proved=annex:p8.pn.six-completeness}{The selected E/EM/EC/EN/EMC/K source presentations}
For each row below and every formula set \(\Gamma\),
\(\Gamma\vdash_L\varphi\) exactly when
\(\Gamma\models^{\ell}_{\mathscr F_L}\varphi\).
Every L-consistent set is satisfiable in its indicated canonical model.
\begin{center}\small
\begin{tabular}{@{}lll@{}}
\toprule
Logic & Frame condition & Canonical choice \\
\midrule
E & arbitrary neighborhoods & minimal or maximal \\
EC & binary intersections & minimal \\
EN & contains the unit & minimal \\
EM & supplemented & supplemented minimal \\
EMC & supplemented and binary intersections & supplemented minimal \\
EMCN = K & filters, not necessarily proper & supplemented minimal \\
\bottomrule
\end{tabular}
\end{center}
All frames and satisfying models are nonempty. These are strong local
completeness claims, not merely empty-premise validity or a finite-model
decision theorem \cite[thms.~2.38--2.42, ex.~2.20]{Pacuit2017}.
\end{hbtheorem}

\subsection{Filters, augmentation and the relational bridge}

\begin{hbtheorem}{thm:p08-pn-augmentation}{proved=annex:p8.pn.augmentation}{Principal neighborhoods and strong completeness on augmented frames}
For a relation R on nonempty W, define
\(N^R(w)=\{X\subseteq W:R(w)\subseteq X\}\). This family
is augmented. Conversely, for augmented N define
\(wR_Nv\) when \(v\in\bigcap N(w)\); then
\(N^{R_N}(w)=N(w)\). These conversions preserve formula truth
at every world with the same valuation, and frame validity under
every valuation. K is strongly complete on augmented neighborhood
frames as well as on all filter frames.

The two frame classes differ: the cofinite filter on \(\N\)
does not contain its empty core and is not augmented. An infinite
filter's existence alone does not supply an equivalent principal
neighborhood relation on that same world set
\cite[\S2.2.1, ex.~2.21]{Pacuit2017}.
\end{hbtheorem}

\begin{hbdefinition}{def:p08-pn-impossible-worlds}{The source's relational semantics for regular logics}
A relational model with impossible worlds is
\(\langle W,W^N,R,V\rangle\), where W is nonempty,
\(W^N\subseteq W\), \(R\subseteq W\times W\), and
\(V:\mathsf{At}\to\wp(W)\). Boolean evaluation is classical
at \emph{every} world. Its modal clause is
\[
 \mathcal M,w\models\Box\varphi
 \quad\Longleftrightarrow\quad
 w\in W^N\text{ and }R(w)\subseteq[![\varphi]!]_{\mathcal M}.
\]
Thus at an impossible world, outside \(W^N\), every boxed
formula is false and every diamond formula is true. Impossible
here means this specified failure of modal necessity, not an
inconsistent Boolean valuation. Its equivalent neighborhoods
are \(N(w)=\{X:R(w)\subseteq X\}\) for \(w\in W^N\)
and \(N(w)=\varnothing\) otherwise
\cite[\S2.2.4, def.~2.22, ex.~2.13]{Pacuit2017}.
\end{hbdefinition}

\begin{hbtheorem}{thm:p08-pn-regular-relational}{proved=annex:p8.pn.regular-relational}{The regular relational boundary and its full canonical construction}
The displayed translation preserves every formula's truth.
Its families are regular; they need not contain the unit.
EMC is sound and strongly locally complete for relational
models with impossible worlds. The proof uses maximal EMC
theories as worlds, \(W^N=\{\Gamma:\Box\topv\in\Gamma\}\),
and \(\Gamma R\Delta\) exactly when every \(\alpha\) with
\(\Box\alpha\in\Gamma\) belongs to \(\Delta\).
This is a separately justified regular completeness construction,
not a claim that every infinite regular neighborhood family is
principal or pointwise representable by one such relation.
\end{hbtheorem}

\subsection{Finite constructions and the actual decidability branch}

\begin{hbdefinition}{def:p08-pn-filtration}{Source quotient and finite generating families}
Expand defined connectives before forming a finite subformula-closed
set \(\Sigma\). Two worlds agree under \(\sim_\Sigma\) when
they agree on every member. Let \(q(w)=[w]_\Sigma\) and
\([X]_\Sigma=q[X]\). The finest generating family is
\[
 \mathcal B([w])=
 \{[![![\psi]!]_{\mathcal M}]_\Sigma:
                 \Box\psi\in\Sigma, [![\psi]!]_{\mathcal M}\in N(w)\}.
\]
The nested brackets mean the quotient image of the original truth set.
The source filtration condition for every tracked \(\Box\psi\)
is exact membership equivalence between that original truth set
in \(N(w)\) and its quotient image in \(N^f([w])\).
Use \(V^f(p)=q[V(p)]\) for atoms in \(\Sigma\); other atoms
may be assigned arbitrarily. For the selected frame properties,
close \(\mathcal B\) under nonempty finite intersections when
C is present, supplement when M is present, and add the quotient
unit when N is present. A unit need not be generated by an
otherwise empty family.
\end{hbdefinition}

\begin{hbtheorem}{thm:p08-pn-filtration}{proved=annex:p8.pn.filtration}{All tracked truth and the simultaneous closure cases}
The quotient constructions just defined preserve truth of every
\(\Sigma\)-formula, have at most \(2^{\#\Sigma}\) worlds,
and retain the requested M/C/N properties of the original model.
For an original proper filter, the filter construction still
excludes the empty proposition. Neither arbitrary quotient
images nor unconstrained closure operations are asserted to
preserve truth: the proof checks every newly added tracked
truth set against the original neighborhood family
\cite[thm.~2.47, lemma~2.17, ex.~2.23]{Pacuit2017}.
\end{hbtheorem}

\begin{hbdefinition}{def:p08-pn-finite-canonical}{Consistent descriptions and source labels}
For \(\Sigma_\varphi=\operatorname{Sub}(\varphi)\), a
description is \(D=X\cup\{\neg\sigma:
\sigma\in\Sigma_\varphi\setminus X\}\), with
\(X\subseteq\Sigma_\varphi\). Retain only L-consistent
descriptions, and give each its specified Lindenbaum extension
\(\Gamma_D\). Put \(W_\varphi=\{\Gamma_D\}\).
A fixed-order conjunction of D is its label \(\lambda_D\);
for \(Y\subseteq W_\varphi\) let \(\lambda_Y\) be the
disjunction of the corresponding labels, using
\(\lambda_\varnothing=\varphi\land\neg\varphi\).
Define \(|\psi|_\varphi=\{\Gamma_D:\psi\in\Gamma_D\}\).
The source finite canonical model has
\[
 N_\varphi(\Gamma_D)=\{Y\subseteq W_\varphi:
                         \Box\lambda_Y\in\Gamma_D\},\qquad
 V_\varphi(p)=|p|_\varphi.
\]
This is an existence construction using classical consistency,
not an algorithm that decides consistency before the decidability
theorem has been proved. Inconsistent descriptions have no
maximal extension and are not assigned a representative.
\end{hbdefinition}

\begin{hbtheorem}{thm:p08-pn-finite-canonical}{proved=annex:p8.pn.finite-canonical}{The complete Boolean-label and non-iterative frame argument}
For Boolean combinations \(\psi\) of \(\Sigma_\varphi\),
\(\vdash_L\psi\) exactly when every \(\Gamma_D\) contains
\(\psi\). Each Y is precisely \(|\lambda_Y|_\varphi\).
The finite canonical model agrees with membership on
\(\Sigma_\varphi\), its Boolean closure, and boxes of that
Boolean closure. For a logic axiomatized over E by non-iterative
schemas, this finite frame validates all its axioms under every
valuation. Every non-theorem \(\varphi\) is falsified on one
such frame with at most \(2^{\#\Sigma_\varphi}\) worlds.
The proof uses L-consistency, not a mistaken identification
of every L-theorem with a propositional tautology.
\end{hbtheorem}

\begin{hbmethod}{def:p08-pn-finite-search}{A finite generator search for the eight M/C/N combinations}
Given \(\varphi\), expand its abbreviations, compute
\(s=\#\operatorname{Sub}(\varphi)\), and consider
\(1\leq j\leq2^s\) labeled worlds. At every world enumerate
a list of at most s arbitrary generating subsets, apply exactly
the M/C/N closure operations of the chosen logic, and enumerate
valuations of the input's occurring atoms. Evaluate the formula
bottom-up using exact neighborhood membership. A satisfying
candidate witnesses satisfiability; if the finite search ends
without one, the formula is unsatisfiable in the selected class.
The same method tests theoremhood by searching for a model of
\(\neg\varphi\). It includes EMN and ECN as well as the six
main comparison rows. No consistency oracle is part of the search.
\end{hbmethod}

\begin{hbtheorem}{thm:p08-pn-decidability}{proved=annex:p8.pn.decidability}{Termination, completeness and the encoding boundary}
The generator search terminates and decides satisfiability and
theoremhood for all eight fixed M/C/N combinations. A model's
bounded generator description has at most \(s j^2\) membership
bits before closure, and valuations require at most s j bits.
This supports a double-exponential exhaustive search in s;
it is not a sharp complexity classification. There are at most
linearly many distinct subformulas in an expanded formula, while
the number of descriptions can be exponential.

For a general non-iterative extension, the finite canonical
construction proves finite-model completeness. Turning that
argument into an effective decision requires a finite,
effective list of non-iterative axiom schemas over E, whose
validity on each finite candidate is checkable under all valuations.
It is not asserted for an
arbitrary nonrecursive infinite set of axioms. The source's
cited NP/PSPACE results are separate sharper results, not
proved by this enumeration \cite[\S\S2.4.1--2.4.2]{Pacuit2017}.
\end{hbtheorem}

\subsection{Extensional structural and general-frame boundaries}

\begin{hbdefinition}{def:p08-pn-bounded-morphisms}{Source bounded morphisms and disjoint unions}
A neighborhood bounded morphism \(f:\mathcal M\to\mathcal M'\)
is a total map of world sets preserving atomic valuation and
satisfying \(f^{-1}(Y)\in N(w)\) exactly when
\(Y\in N'(f(w))\), for every subset Y of the target world set.
For models on pairwise disjoint domains, their disjoint union
has union valuation and, at \(w\in W_i\),
\(Y\in N(w)\) exactly when \(Y\cap W_i\in N_i(w)\).
This is an extensional inverse-image condition, not a monotonic
bisimulation condition silently applied to nonmonotonic models.
\end{hbdefinition}

\begin{hbtheorem}{thm:p08-pn-structural-truth}{proved=annex:p8.pn.structural-truth}{Full morphism and disjoint-component truth preservation}
Bounded morphisms preserve every formula's truth at corresponding
worlds; surjectivity is not needed for that statement. Component
embeddings into the source disjoint union are bounded morphisms.
Thus adding other components does not change truth at an old
world. Supplementation and binary intersection/unit properties
are retained by the disjoint-union construction when present
in every component \cite[props.~2.3--2.5]{Pacuit2017}.
\end{hbtheorem}

\begin{hbtheorem}{thm:p08-pn-one-step}{proved=annex:p8.pn.one-step}{The source one-step E and EM characterizations with theorem-level premises}
For finite lists of atoms \(\vec p,\vec q\) and arbitrary
formula lists \(\varphi_1,\ldots,\varphi_k\),
\(\psi_1,\ldots,\psi_m\), the sequent
\[
 \vec p,\Box\varphi_1,\ldots,\Box\varphi_k
 \ \Rightarrow\ \vec q,\Box\psi_1,\ldots,\Box\psi_m
\]
is valid on monotonic frames exactly when the atomic lists
overlap or some \(\varphi_i\to\psi_j\) is globally valid
on that class. On arbitrary neighborhoods the second condition
is instead global validity of \(\varphi_i\leftrightarrow\psi_j\).
These are theorem-level matrix conditions. An implication
derived from arbitrary unchanged local contexts cannot be
boxed on both sides. For EC the source's conjunction-folding
step is sound from an already valid sequent with a boxed
conjunction to one with its two separately boxed antecedents;
it supplies no monotonic projection absent from EC. More exactly,
on intersection-closed frames the displayed sequent is valid
exactly when the atomic lists overlap or there is a nonempty
subset \(I\subseteq\{1,\ldots,k\}\) and a j such that
\(\bigwedge_{i\in I}\varphi_i\leftrightarrow\psi_j\)
is valid throughout that class. This includes a singleton I;
empty intersections would incorrectly add N to EC.
\end{hbtheorem}

\begin{hbdefinition}{def:p08-pn-general-frames}{Admissible proposition sets}
A general neighborhood frame is \(\langle W,N,\mathcal A\rangle\),
where \(\mathcal A\subseteq\wp(W)\) is a Boolean algebra
containing the unit and closed under
\(m_N(X)=\{w:X\in N(w)\}\). Valuations must take their
atomic values in \(\mathcal A\). For a consistent classical
logic L, the source canonical general frame uses
\(W=M_L\), \(N=N_L^{\min}\), and \(\mathcal A=P_L\).
Uniform formula-substitution closure of L is an explicit
hypothesis for its all-admissible-valuations validity claim.
\end{hbdefinition}

\begin{hbtheorem}{thm:p08-pn-general-completeness}{proved=annex:p8.pn.general-completeness}{Canonical model truth versus admissible frame validity}
For every consistent substitution-closed classical L, its
canonical general frame validates L under every admissible
valuation. L is sound and strongly locally complete for the
class of its general neighborhood frames. This restriction
to admissible propositions is material: it is not unrestricted
ordinary-frame completeness for every classical extension
\cite[thm.~2.45, cor.~2.3]{Pacuit2017}.
\end{hbtheorem}

\paragraph{Attribution and scope.}
The canonical extrema, proof-set membership tests, supplemented EM
construction and combined filter application here follow the published
2017 \emph{Core Theory} presentation. The earlier Chellas treatment in
Section~\ref{sec:p08-n-b-006-03} uses its own source notation and
closure constructions; the shared sheets agree after the explicit
defined-diamond and theorem-only local-deduction bridges above.
Specifically Chellas's least monotonic M is Pacuit's EM, its
least regular R is Pacuit's EMC (also ECM), and its normal K is
Pacuit's EMCN: each pair has the same translated propositional
base, schemas and theorem rules. This assertion concerns the
specified nonempty model classes; Chellas's convention permitting
empty domains does not provide a world satisfying premises and
does not change these theorem or local-consequence comparisons.
Scott's and Montague's original presentations remain separately
attributed. No earliest source is credited with all these later
constructions, and EN is never identified with K.
